\documentclass{article}
\usepackage[UTF8]{ctex}
\usepackage{amsmath,amssymb,amsfonts}
\usepackage{enumitem}
\usepackage{geometry}
\geometry{a4paper, margin=1in}
% === 错题本排版优化（不参与题目标识） ===
\setlength{\parindent}{0pt}
\setlength{\parskip}{0.35em plus 0.10em minus 0.08em}
\setlength{\emergencystretch}{3em}
\setlist{itemsep=0.25em, topsep=0.40em, parsep=0pt}
\allowdisplaybreaks[2]
% === 错题本排版优化结束 ===
\begin{document}
\section*{中山大学2026年数学分析真题}
\section*{154．中山大学2026年数学分析试卷}
\begin{itemize}

\subsection{题目1}
  \item[1．] 求极限 \(\lim _{x \rightarrow 0^{+}}\left(\frac{\sin x}{x}\right)^{\frac{1}{x^{2}}}=\) \(\_\_\_\_\)．

\subsection{题目2}
  \item[2．] 设函数 \(y=y(x)\) 由方程 \(\ln (x y)=x^{2}-y^{2}\) 隐式定义，求点 \((x, y)=(1,1)\) 处的导数 \(\frac{\mathrm{d} y}{\mathrm{~d} x}(1,1)=\)\\
\(\_\_\_\_\)．
\end{itemize}

\subsection{题目3}
3．求定积分 \(\int_{0}^{\pi} \frac{\mathrm{d} x}{\sqrt{2}+\sin x}=\) \(\_\_\_\_\)．

\subsection{题目4}
4．曲线由参数方程 \(x=t, y=t^{2}\) 给出，求该曲线在 \(t=1\) 处的曲率 \(\_\_\_\_\)．

\subsection{题目5}
5．已知函数 \(f(x, y)=\left\{\begin{array}{ll}\frac{x y}{\left(x^{2}+y^{2}\right)^{p}}, & x^{2}+y^{2} \neq 0 \\ 0, & x^{2}+y^{2}=0\end{array}\right.\) 在 \((0,0)\) 处连续，求 \(p\) 的取值范围 \(\_\_\_\_\)．\\
\(\_\_\_\_\)

\subsection{题目6}
6．求函数 \(\mu(x, y, z)=x y^{2}+z^{3}-x y z\) 在点 \((1,1,2)\) 处沿方向 \(\mathbf{l}=\left(\frac{1}{2}, \frac{\sqrt{2}}{2}, \frac{1}{2}\right)\) 的方向导数 \(\_\_\_\_\)．

\subsection{题目7}
7．计算积分 \(\iint_{0 \leq s \leq 1, s \leq t \leq \sqrt{s}} \frac{\sin t}{t} \mathrm{~d} t \mathrm{~d} s=\)

\subsection{题目8}
8．设 \(L\) 为球面 \(x^{2}+y^{2}+z^{2}=4\) 被平面 \(x+y+z=0\) 所截得的圆周，计算积分 \(\int_{L} x^{2} \mathrm{~d} s=\)

\subsection{题目9}
9．级数 \(\sum_{n=2}^{\infty} \frac{1^{n}+2^{n}+\cdots+10^{n}}{n(\ln n)^{2}} e^{n x}\) 的一致收敛区域为 \(\_\_\_\_\)．

\subsection{题目10}
10．如果级数 \(\sum_{n=1}^{\infty}\left(1-\frac{p \ln n}{2 n}\right)^{n}\) 是收敛的，则 \(p\) 的取值范围是 \(\_\_\_\_\)．

\subsection{题目11}
11．设 \(S(x)=\sum_{n=1}^{\infty} \frac{\sin 2 n x}{n}\) ，则 \(S\left(\frac{\pi}{6}\right)=\)

\subsection{题目12}
12． \(\int_{0}^{1} \frac{1}{x} \ln \frac{1+x}{1-x} \mathrm{~d} x=\)\\

\subsection{题目2}
原题拆分片段：
二（15 分）设 \(f(x)\) 在 \((0,1)\) 上有定义，\(f(x) \tan x\) 单调递增， \(\tan (f(x))\) 单调递减，证明：\(f(x)\) 在 \((0,1)\)上连续。

三．（15 分）证明：函数 \(f(x)=x \cos \frac{1}{x}\) 在 \((0,+\infty)\) 上一致连续．\\

\section{原题拆分}
\subsection{题目1302}
原题拆分片段：
且（15 分）计算积分 \(I(b)=\int_{0}^{\frac{\pi}{2}} \frac{\mathrm{~d} x}{b+\tan x}, b>0\) ．

五．（15 分）设 \(f(x)\) 在 \([0,+\infty)\) 上一致连续，且对任意固定的 \(x \in[0,1]\) ，有 \(\lim _{n \rightarrow \infty} f(x+n)=0\) ，证明：函数列 \(\{f(x+n) \mid n \geq 1\}\) 在 \([0,1]\) 上一致收敛于 0 。

六．（15分）设函数列 \(\left\{a_{n}(x)\right\}\) 在 \([0,1]\) 上有定义，且存在不依赖于 \(x, n\) 的常数 \(K\) ，使得

\[
\left|a_{1}(x)\right| \leq K, \sum_{n=1}^{m}\left|a_{n}(x)-a_{n+1}(x)\right| \leq K, \forall m \geq 1 .
\]

证明：如果级数 \(\sum_{n=1}^{\infty} c_{n}\) 收敛，\(c_{n} \geq 0\) ，则级数 \(\sum_{n=1}^{\infty} c_{n} a_{n}(x)\) 在 \([0,1]\) 上一致收敛。

七．（15 分）令 \(B_{R}=\left\{\left(x_{1}, x_{2}\right) \in \mathbb{R}^{2} \mid x_{1}^{2}+x_{2}^{2} \leq R^{2}\right\}\) 是 \(\mathbb{R}^{2}\) 中的半径为 \(R\) 的闭球，\(R>0\) ．设 \(f(x)\) 是\(B_{2}\) 上的连续函数，满足条件：对任意的 \(\mathbf{x}, \mathbf{y} \in B_{2}\) ，成立 \(\frac{f(\mathbf{x})+f(\mathbf{y})}{2} \geq f\left(\frac{\mathbf{x}+\mathbf{y}}{2}\right)\) ．证明：\\
（1）对任意的 \(\mathbf{x}, \mathbf{y} \in B_{2}\) 和任意的 \(\lambda \in(0,1)\) ，成立 \(\lambda f(\mathbf{x})+(1-\lambda) f(\mathbf{y}) \geq f(\lambda \mathbf{x}+(1-\lambda) \mathbf{y})\) ．\\
（2）存在一个常数 \(C>0\) ，使得如下不等式成立：

\[
|f(\mathbf{x})-f(\mathbf{y})| \leq C\|\mathbf{x}-\mathbf{y}\|, \forall \mathbf{x}, \mathbf{y} \in B_{1} .
\]

这里 \(\|\mathbf{x}-\mathbf{y}\|\) 表示 \(\mathbf{x}=\left(x_{1}, x_{2}\right)\) 和 \(\mathbf{y}=\left(y_{1}, y_{2}\right)\) 之间的距离 \(\sqrt{\left(x_{1}-y_{1}\right)^{2}+\left(x_{2}-y_{2}\right)^{2}}\) ．
\end{document}
